\section{DL dataset and regularization}
\label{sec:regularization}

\textbf{What regularizes the graph models.} Four mechanisms act at once, all fixed before any
test label is read: dropout $0.2$ after every message-passing block and inside the head; weight
decay $10^{-5}$ in Adam; early stopping on validation mean ROC-AUC with patience 25, which caps
the effective number of epochs (the frozen models stop between epoch 30 and 61, far short of the
200-epoch budget); and the masked multi-task loss itself, which makes the twelve endpoints share
one encoder and so acts as the inductive bias that a single sparse endpoint cannot overfit alone.
Batch normalisation between layers is used for optimisation stability rather than regularization.

\textbf{Augmentation study.} The plan asks whether graph augmentation adds anything on top. We
compare the frozen GCN against two augmentations applied to \emph{training batches only}, so
validation and test molecules, early stopping and every reported metric describe unperturbed
structures:

\begin{itemize}
  \item \texttt{gcn-augmask} --- \textbf{random atom masking}, 10\%: the feature vector of a
        randomly chosen 10\% of atoms is zeroed. The atom and its bonds stay in the graph, so the
        topology, the batch vector and the pooling denominators are unchanged; only the atom's
        identity is hidden, and the model must infer it from its neighbourhood.
  \item \texttt{gcn-augdrop} --- \textbf{random bond dropping}, 10\%: 10\% of bonds are deleted
        with their features. Both directions of a bond are removed together, so the graph stays
        undirected as the message passing assumes.
\end{itemize}

Both are re-sampled every epoch, drawn from the checkpointed RNG so a resumed run replays the
same augmentations. Everything else --- architecture, frozen hyper-parameters, loss, optimiser,
early stopping, seeds --- is identical to the baseline, which is the un-augmented \texttt{gcn}
run of Table~\ref{tab:gnn-main}; the three arms therefore differ in exactly one factor.

\begin{table}[t]
\centering
\caption{Graph augmentation on \texttt{scaffold\_det}: test ROC-AUC and PR-AUC, 12-endpoint
mean, mean $\pm$ std over seeds 0--2. The baseline is the un-augmented GCN. The mean test
positive rate is 0.103, which is the PR-AUC of a random classifier.}
\label{tab:augmentation}
\begin{tabular}{lcc}
\hline
Arm & ROC-AUC & PR-AUC \\
\hline
\texttt{gcn} (no augmentation) & $0.752 \pm 0.008$ & $0.354 \pm 0.004$ \\
\texttt{gcn-augmask} (atom masking 10\%) & $0.749 \pm 0.005$ & $0.358 \pm 0.005$ \\
\texttt{gcn-augdrop} (bond dropping 10\%) & $0.753 \pm 0.005$ & $0.340 \pm 0.022$ \\
\hline
\end{tabular}
\end{table}

\textbf{Result: neither augmentation helps.} Every difference in Table~\ref{tab:augmentation} is
smaller than the seed-to-seed spread it sits next to. Atom masking costs $0.003$ ROC-AUC and
gains $0.004$ PR-AUC; bond dropping gains $0.001$ ROC-AUC and loses $0.014$ PR-AUC, with a
PR-AUC standard deviation ($0.022$) five times the baseline's, i.e.\ it mostly adds variance.
With three seeds and differences this small we claim no ordering between the three arms: the
honest reading is that 10\% augmentation neither helps nor hurts this model on this split.

That is a coherent result rather than a disappointing one. The GCN is already stopped early
(epoch 30--61 of 200) by a validation signal computed on 12 endpoints at once, and its capacity
is modest; the evidence that it is not over-fitting is the MPNN ablation grid in
Table~\ref{tab:gnn-ablation}, where every arm but one --- including the deliberately weakened
ones, such as bond features removed or positive weights switched off --- scores at or above the
frozen baseline (the exception is the 5-layer arm, $0.731$ against $0.732$, a gap far inside its
own spread). A model that improves when it is weakened is limited by its tuning budget, not by
excess capacity. Augmentation addresses over-fitting, and over-fitting is not what limits this
model, so the absence of a gain is what the rest of our evidence predicts. The regularization
that does matter here is the multi-task masked loss, which is tested directly by the
multi-task-versus-single-task ablation.
